\documentclass[conference,10pt]{IEEEtran}
\usepackage{times,graphicx,amsmath,booktabs,url,balance}
\usepackage[hidelinks]{hyperref}
\renewcommand{\arraystretch}{1.15}
\begin{document}

\title{fpgAI: Agent-Designed FPGA Scaleout for Local LLM Inference}
\author{Group X\\
Sameer Suleman, Yax Patel, Vaibhav [surname], Siddh [surname]\\
Supervisor: Dr. [Supervisor]}
\maketitle

\section{Executive Summary}
Running a large language model locally is limited less by arithmetic than by the
hardware available to run it on. Purpose-built accelerators cost tens of thousands
of dollars and take years to design, while the programmable hardware a student or
small lab can afford is individually too small to hold a useful model. The design
process itself is the bottleneck: a chip is specified years before the models it
will run, and closing that loop currently takes teams of specialists months.

We are building \textbf{fpgAI}, a system in which software agents take a language
model's specification as input and generate the hardware required to run it. The
agents derive the compute datapath and the inter-board communication fabric from
the model's own properties, write the RTL, carry it through simulation, synthesis,
static timing analysis, FPGA technology mapping and place-and-route without a human
in the loop, and emit a measured performance profile. A sizing layer then determines
how many FPGA boards the model needs and validates that answer packet by packet
before any hardware is purchased. Capacity grows by adding boards rather than by
buying a larger one.

The end users are students, hobbyists, university research groups and small
hardware teams who need local, private inference on hardware they can actually
afford, and researchers studying automated chip design. Our target is
Qwen3-0.6B running at 40 or more tokens per second across two FPGA boards costing
roughly \$300 each. For comparison, Architect Labs measured 12.1 tokens per second
on the same model using a \$17{,}995 AMD Versal board~\cite{redwood}.

The principal risks are the runtime cost of place-and-route, which is roughly
500 times slower than the rest of the loop and therefore constrains how agents may
iterate, the unproven ability of a language model to correct its own RTL from tool
feedback, and the platform dependency that vendor place-and-route tools do not run
on all team machines. Each is addressed in Sections~\ref{sec:modules}
and~\ref{sec:monitor}. The software half of the system is already implemented and
covered by 66 automated tests, so the project begins from a working base rather
than from zero.

\section{Expression of Need}
Verifying and deploying a new hardware design is slow and specialist work. A
conventional flow runs from specification to RTL, through verification, synthesis
and timing closure, and finally through physical implementation, with human
engineers interpreting tool output and revising the design at every stage. For
inference accelerators the problem compounds, because the chip must be specified
years before the models it will eventually run, and by the time silicon exists the
models have changed.

At the other end of the scale, a student or small lab that wants to run a model
locally faces a different wall. A single affordable FPGA board holds only a few
hundred kilobytes of fast on-chip memory and a gigabyte of slower external memory,
which is not enough for a useful model on its own. Commercial accelerators that are
large enough cost more than the entire budget of a research group, and cloud
inference means sending data off-premises, which is unacceptable for medical,
legal and industrial users.

The stakeholders are therefore twofold. First, users who need affordable, private,
local inference: university labs, hardware startups, and anyone handling data that
cannot leave the building. Second, the hardware design community, for whom
automating the path from specification to working silicon shortens a cycle
currently measured in engineer-months.

What is missing is a system that closes both gaps at once: one that derives the
hardware from the model instead of guessing years ahead, and that reaches useful
capacity by combining cheap boards instead of requiring one expensive one.

\section{Background}
\textbf{Redwood} (Architect Labs, 2026) is the closest prior work. An autonomous AI
system produced the RTL, UVM verification environment, formal properties, firmware
and drivers for an inference accelerator in under two weeks, reaching over 95\%
functional coverage, and deployed it on an AMD Versal VPK180 running Qwen3-0.6B at
12.1 tokens per second~\cite{redwood}. Two things it did not do define our
contribution: the architecture was specified by two human architects rather than
derived from the model, and the result ran on a single chip with no inter-chip
scaleout. Notably, its measured throughput was below that of a \$250 NVIDIA Jetson
Orin Nano, confirming that the result was about the design process rather than
performance.

\textbf{MAGE}~\cite{mage} demonstrated that multi-agent LLM systems reach 95.7\% on
the VerilogEval-Human benchmark, establishing that the underlying RTL generation
capability exists. \textbf{AlphaChip}~\cite{alphachip} applied reinforcement
learning to floorplanning, and \textbf{ChipNeMo}~\cite{chipnemo} built a
domain-adapted LLM for chip designers. Commercially, Synopsys and Cadence both ship
AI-driven optimization, but all of these automate a single stage rather than the
path from specification to bitstream.

On the scaleout side, tensor parallelism with ring all-reduce collectives is
standard practice in GPU clusters~\cite{megatron}, and Ethernet-native multi-chip
fabrics are an active commercial direction. No existing work generates that fabric
as part of the same agentic loop that generates the compute.

\section{Objectives}
Objectives are quantifiable and measured on hardware unless stated otherwise.

\begin{table}[h]
\centering
\caption{Technical objectives}
\begin{tabular}{@{}p{4.9cm}p{2.9cm}@{}}
\toprule
\textbf{Metric} & \textbf{Target} \\
\midrule
Distinct hardware blocks generated and signed off autonomously & $\geq 4$ \\
Automated signoff gates in the loop & 5 (simulate, synthesize, time, map, place-and-route) \\
Specs reaching full signoff within 5 agent iterations & $\geq 90\%$ \\
Blocks corrected by an LLM from parsed tool feedback alone & $\geq 1$, transcript captured \\
Qwen3-0.6B throughput, two boards & $\geq 40$ tokens/s \\
Analytic prediction versus measured throughput & within $15\%$ \\
Numerical agreement across 1 to 8 boards & bit-exact, $100\%$ \\
Inter-board link rate & $\geq 1$ Gbps sustained \\
Data integrity at bit error rate $10^{-6}$ & zero corrupted results \\
Automated regression tests & $\geq 100$ \\
Functional coverage on generated blocks & $\geq 90\%$ \\
Total component budget & $\leq$ \$800 CAD \\
\bottomrule
\end{tabular}
\end{table}

\section{Technical Modules}\label{sec:modules}
The system divides into four modules, one owned by each group member. They are
separable because the interfaces between them are machine-readable files that
already exist: a JSON performance profile emitted by Module~A and consumed by
Module~B, a constraints-and-wrapper contract between Modules~A and~C, and a
testbench contract between Modules~A and~D. Each module can therefore be developed
and validated against a stub of its neighbours.

\subsection{Module A: Agentic Generation and Signoff (Sameer)}
The orchestrator and the agents. Derives block specifications from the model
(datapath width from the quantization, accumulator width from the deepest
reduction so that overflow is impossible by construction), prompts an LLM for RTL,
parses tool failures into structured feedback, and iterates to signoff. Also owns
the sizing layer that models compute, memory bandwidth, on-chip residency, KV cache
traffic and network cost. Validated in simulation, so it needs no hardware.

\subsection{Module B: Interconnect Fabric and Collectives (Yax)}
The inter-board communication layer: framing, CRC, sequence numbers, go-back-N
retransmission, credit-based flow control, and ring all-reduce for arbitrary board
counts. Also owns the multi-agent architecture, defining the writer, verifier and
reviewer roles and how tool output is routed between them. Entirely software, so
it is never blocked waiting for hardware.

\subsection{Module C: Physical Implementation (Vaibhav)}
Everything that runs on real silicon: the Vivado project, constraints, top-level
wrapper, place-and-route, bitstream generation and board bring-up. Feeds real
utilization and post-route timing back into the profile, replacing estimates. Uses
the board already in hand, so it starts immediately.

\subsection{Module D: Verification and Physical Link (Siddh)}
Raising verification from golden-model checks to coverage-driven and formal
methods, including a proof that the derived accumulator cannot overflow for any
model. Also owns the physical board-to-board link: signalling, cabling, achievable
data rate, and the bring-up and acceptance test plan. The verification half is
software and the link analysis is paper work, so neither waits on hardware.

\section{Deliverables}
Three tiers. \textbf{Bronze is the minimum viable product} and is deliberately
scoped so that it depends on neither a hardware purchase nor the unproven ability
of an LLM to correct its own RTL, which are the two riskiest items in the project.
\textbf{Silver is the core result} the project should be judged on. \textbf{Gold is
stretch work}, pursued only once Silver is demonstrated and dropped first if the
schedule slips.

\subsection{Bronze: minimum viable product}
\emph{Modules A, C and D. No purchase needed; both boards are already in hand.}
\begin{itemize}\itemsep1pt \parskip0pt
\item Agents generate and sign off the compute datapath and the fabric endpoint
through all five gates with no human intervention.
\item A bitstream built from agent-written RTL loads on one FPGA and passes its
self-test, closing the specification-to-silicon path for a single block.
\item The simulated cluster hosts Qwen3-0.6B and reports throughput within 15\% of
the analytic prediction.
\item At least 100 automated regression tests pass from one command.
\end{itemize}

\subsection{Silver: core result}
\emph{Adds Modules B and D. Needs the second board, already owned.}
\begin{itemize}\itemsep1pt \parskip0pt
\item Two boards linked over the generated fabric, with framing, checksums and
flow control verified on real hardware.
\item Ring all-reduce across both boards producing results bit-identical to a
single board, including under injected bit errors at $10^{-6}$.
\item Qwen3-0.6B running end to end at 40 or more tokens per second, measured and
compared against the prediction.
\item At least one block demonstrably corrected by an LLM reading tool output,
with the transcript published.
\end{itemize}

\subsection{Gold: stretch}
\emph{All modules. Pursued only after Silver; first to be cut.}
\begin{itemize}\itemsep1pt \parskip0pt
\item Larger models hosted by adding boards, showing throughput holds roughly
constant as model and cluster grow together.
\item A multi-agent swarm in which an agent writes the verification as well as the
RTL, with coverage-driven and formal signoff.
\item A quantified comparison of interconnect options, including photonic links.
\item Optionally, the compute block hardened to GDSII on an open process and
submitted to a shared shuttle.
\end{itemize}

\section{Resources}
\subsection{People}
Four members, one technical module each, as above. Sameer and Yax work primarily in
software; Vaibhav owns the physical flow; Siddh spans verification and link
hardware.

\subsection{Software}
All tooling is free. Icarus Verilog and cocotb for simulation, Yosys for synthesis
and FPGA technology mapping, OpenSTA for static timing, SymbiYosys for formal
properties, AMD Vivado (free tier, which covers the Artix-7 and Zynq-7000 devices
we target) for place-and-route and bitstream generation, and PyTorch for model
quantization. Agent inference costs cents per converged run.

\subsection{Components and budget}
\begin{table}[h]
\centering
\caption{Budget, CAD}
\begin{tabular}{@{}p{4.6cm}rr@{}}
\toprule
\textbf{Item} & \textbf{Qty} & \textbf{Cost} \\
\midrule
Zynq-7000 class FPGA boards (already owned) & 2 & 0 \\
Additional FPGA board & 1 & 520 \\
Ethernet and ribbon cabling, connectors & -- & 60 \\
Pmod adapters and passives & -- & 40 \\
Contingency & -- & 150 \\
\midrule
\textbf{Total} & & \textbf{770} \\
\bottomrule
\end{tabular}
\end{table}

Two boards are already in hand, so the Bronze and Silver deliverables are reachable
with no purchase. Sponsorship from an industry partner is being pursued to cover
the third board and extend the Gold objectives.

\section{Schedule and Milestones}
Milestones fall on two-week boundaries to align with the bi-weekly instructor
meetings. Modules run in parallel throughout.

\begin{table}[h]
\centering
\caption{Bi-weekly milestones}
\begin{tabular}{@{}cp{6.0cm}@{}}
\toprule
\textbf{Wk} & \textbf{Milestone} \\
\midrule
2 & Proposal delivered. Vivado toolchain installed and a reference bitstream builds on the board in hand. \\
4 & Place-and-route added as the fifth signoff gate; utilization and post-route timing reach the profile. \\
6 & Agent-written RTL placed, routed and loaded onto one board, self-test passing. \textbf{Bronze.} \\
8 & LLM iteration demonstrated: a block fixed from tool feedback, transcript captured. Multi-role agent split implemented. \\
10 & Physical board-to-board link characterized; measured data rate feeds the sizing model. \\
12 & Two boards linked; framing and CRC verified on real hardware. \\
14 & Ring all-reduce across two boards, bit-exact against a single board. \\
16 & Qwen3-0.6B hosted end to end on two boards; measured against prediction. \textbf{Silver.} \\
18 & Coverage instrumentation and formal accumulator proof complete. \\
20 & Larger model hosted by adding a board; scaling behaviour characterized. \\
22 & Agent-written testbenches; interconnect comparison study. \textbf{Gold.} \\
24 & Final integration, full regression, report and demonstration. \\
\bottomrule
\end{tabular}
\end{table}

Roughly four weeks of float are held across the second half to absorb hardware
delays, which historically dominate schedule risk on projects of this type.

\section{Progress Monitoring and Control}\label{sec:monitor}
Progress is tracked against artifacts rather than self-reported percentages. Every
milestone above has a binary exit criterion that is either demonstrable or not.

Task tracking uses GitHub issues, one per module, with items owned by name and a
pinned overview issue carrying the timeline and open decisions. All code lives in a
public repository with a regression suite that runs from a single command, so
progress is verifiable by running it. Weekly team meetings review demonstrations
rather than status, and the bi-weekly instructor meetings align with the milestone
boundaries.

Three specific control measures address the known risks. First, because
place-and-route is roughly 500 times slower than the other gates, the loop is
two-tiered: agents iterate freely against the fast gates and place-and-route runs
once when they all pass, with agent retries capped so a single block cannot consume
a day. Second, because vendor place-and-route does not run on every team machine,
bitstream generation is assigned to a single owner with a documented environment
rather than assumed to be available to everyone. Third, if the LLM iteration
objective proves unreachable, the deterministic rule-based agent already in the
repository satisfies the loop requirement, and the claim is narrowed honestly in
the report rather than overstated.

If a module falls behind, members on Gold-level work move to unblock the critical
path, and Gold objectives are dropped before Silver, which is dropped before
Bronze.

\section{Broader Considerations}
\textbf{Health and safety.} All hardware is low-voltage and USB or wall-adapter
powered, presenting minimal electrical hazard. Boards are handled on an ESD mat
with wrist straps, since FPGAs are static-sensitive. Soldering, if required, is
done in a ventilated lab space.

\textbf{Sustainability and environment.} The project reuses existing hardware
rather than purchasing new where possible, and two of the boards are already owned.
Local inference on low-power FPGAs draws a few watts against the hundreds of watts
of a datacenter GPU, and the scaleout approach extends the useful life of small
boards that would otherwise be individually inadequate.

\textbf{Security and privacy.} Inference runs entirely on local hardware and no
user data leaves the machine, which is the central advantage for medical, legal and
industrial users for whom cloud inference is not permissible.

\textbf{Ethics.} All code, results and negative findings are published openly.
Performance figures produced by simulation are labelled as predictions and
distinguished from measurements, and limitations are stated explicitly rather than
omitted.

\textbf{Economic.} The system targets a cost reduction of more than an order of
magnitude against comparable published work, with the explicit aim of making local
inference and automated hardware design accessible to groups that cannot afford
commercial tooling.

\textbf{Human factors and aesthetics.} The system is operated through a single
command and a status display showing per-board activity and throughput, so that a
user need not understand the internal flow to use it. Tool output is summarised
into readable tables rather than raw logs.

\textbf{Regulatory and compliance.} All software tools are used within their
licence terms, including the free tier of the vendor toolchain. Model weights are
used in accordance with their published licences. No radio, mains-voltage or
medical-device regulation applies to the hardware built.

\begin{thebibliography}{9}
\bibitem{redwood} Architect Labs, ``Redwood: A Frontier AI Accelerator Designed,
Verified, and Deployed from Scratch in 2 Weeks by AI,'' arXiv:2608.26418, 2026.
\bibitem{mage} Y. Zhao et al., ``MAGE: A Multi-Agent Engine for Automated RTL Code
Generation,'' arXiv:2412.07822, DAC 2025.
\bibitem{alphachip} A. Mirhoseini et al., ``A graph placement methodology for fast
chip design,'' \emph{Nature}, vol. 594, pp. 207--212, 2021.
\bibitem{chipnemo} M. Liu et al., ``ChipNeMo: Domain-Adapted LLMs for Chip Design,''
arXiv:2311.00176, 2023.
\bibitem{megatron} M. Shoeybi et al., ``Megatron-LM: Training Multi-Billion
Parameter Language Models Using Model Parallelism,'' arXiv:1909.08053, 2019.
\bibitem{yosys} C. Wolf, ``Yosys Open SYnthesis Suite,'' \url{https://yosyshq.net/yosys/}.
\end{thebibliography}

\end{document}
